\section{Main theorem}
\label{sec:main}

\subsection{Setting}
We use the notation of Papers~F and~H and recall only what the statements need.
A periodic window is $p=\mathcal N_V\,\Gamma\,h/Z$, with $\mathcal N_V$ a Gaussian of variance $V=\sigma^2/b$ in $v$, $\Gamma=N(0,L^2)$ in $w$, and
\begin{gather}
h-1=e^{-w^2/2s^2}R(t),\quad t=k'v,\\ R(t)=\sum_{n\ge1}q0_n\cos(nt-\varphi_n).
\end{gather}
We write $x=(k'\sigma)^2$, $\tau=x/b$, $\eta=\sup|R|$ and $\ell=\ln(1/\sigma)$.
A \emph{state} is a choice of $(b,L^2,s^2)$ with $0<b\le1$, (P0) $L^2\le1/\sigma$ and (P4) $\sigma^2\le s^2\le L^2/4$.
The conditions (P0)--(P7) are those of Table~5 of~\cite{F}.
Throughout,
\begin{equation}
B=\frac{b}{\sigma^2}=\frac1V,\qquad \Gamma_\ell=\sqrt{\frac{1+10/\ell}{\ell}} .
\end{equation}
As in Eq.~(3) of~\cite{F},
\begin{equation}
g\ \ge\ (1-|\varepsilon|)\,\bigl((a)+(a')\bigr),\quad (a)=\frac{2A_w}{\sigma\,\mathrm{den}},\quad (a')=\frac{N_\phi}{\sigma\,\mathrm{den}},
\label{eq:ga}
\end{equation}
with $\mathrm{den}=\sigma^2(\Jvv+\Jww)$ and $|\varepsilon|\le\bar\varepsilon$ the window-class constant of (P6).
All theorems below bound $(a)+(a')$ from below; the factor $1-\bar\varepsilon$ enters only through Eq.~\eqref{eq:ga}.

\emph{The class.} A window is \emph{admissible} if it satisfies (P0), (P2)--(P4) and (P6), and $5\times10^{11}\sigma^2\le x\le0.003$.
This is the class of Theorem~H1 of~\cite{H}.

\subsection{Statement}

\begin{theorem}[every focus, every resolution] \status{proof, computer-assisted: Theorems~\ref{thm:narrow}, \ref{thm:wide} and Theorem~H1 of~\cite{H}}
\label{thm:main}
Let $\ell\ge55.26$ and let an admissible window satisfy $\eta\le1/\ell$.
Then every state satisfies
\begin{equation}
g\sqrt\ell\ \ge\ 1-\bar\varepsilon .
\end{equation}
If $\ell\ge960.8$, every state satisfies the stronger bound
\begin{equation}
g\sqrt\ell\ \ge\ (1-\bar\varepsilon)\sqrt{1+10/\ell}\,;
\end{equation}
for $55.26\le\ell<960.8$ the stronger bound holds for every state with $B\ge2.5\times10^{10}$.
\end{theorem}

In the observer's words: at every resolution finer than about $10^{-24}$, and for every position and width of the focus, the information that survives the blur is at least the fraction $(1-\bar\varepsilon)/\sqrt\ell$.
There is no state in which a clever focus does better than the exchange rate.

Theorem~\ref{thm:main} is assembled from five statements, each for one kind of focus.
Four are new; the fifth is Theorem~H1 of~\cite{H}.
We state them in the language of windows; the sections that prove them give the sharper technical forms.

\begin{theorem}[narrow focus] \status{proof, computer-assisted: Sec.~\ref{sec:narrow} and App.~\ref{app:C}}
\label{thm:narrow}
Let $\ell\ge55.26$ and let an admissible window satisfy $\eta\le1/\ell$.
Every state with
\begin{equation}
b\le\tfrac12,\qquad B\ge\max\bigl(2.5\times10^{10},\,50(8\ell+100)\bigr)
\end{equation}
satisfies $(a)+(a')\ge\Gamma_\ell$.
There is no upper limit on $\ell$.
\end{theorem}

\begin{theorem}[focus at the resolution] \proofst
\label{thm:res}
Let $\ell\ge55.26$ and let an admissible window satisfy $\eta\le1/\ell$.
Every state with $\tfrac12<b\le1$ satisfies $(a)+(a')\ge0.1995$; if the ripple has energy in the focus ($G>0$ in the notation of Sec.~\ref{sec:setting}) then $(a)+(a')\ge0.6$.
In particular $(a)+(a')\ge\Gamma_\ell$.
\end{theorem}

\begin{theorem}[wide focus] \status{proof, computer-assisted: Sec.~\ref{sec:wide} and App.~\ref{app:C}}
\label{thm:wide}
Let $\ell\ge960.8$ and let an admissible window satisfy $\eta\le1/\ell$.
Every state with $8\ell+100\le B<\ell^4$ satisfies $(a)+(a')\ge\Gamma_\ell$.
\end{theorem}

\begin{theorem}[very wide focus] \proofst
\label{thm:verywide}
Let $\ell\ge960.8$ and let a periodic window satisfy (P0), (P3), (P4) and $5\times10^{11}\sigma^2\le x\le0.003$.
Every state with $0<B\le8\ell+100$ satisfies $(a)+(a')\ge\Gamma_\ell$.
Neither $\eta\le1/\ell$ nor (P2) is used.
\end{theorem}

Theorem~\ref{thm:verywide} follows from the next two statements; the first holds where the focus spans many periods, the second everywhere.
Put $\tau_0(\ell)=4\ln\ell+40$.

\begin{proposition}[the mean-square bound without limits on $B$ and $\ell$] \proofst
\label{prop:BE}
Under the hypotheses of Theorem~\ref{thm:verywide}, every state with $0<B\le8\ell+100$ and $\tau\ge\tau_0(\ell)$ satisfies
\begin{equation}
(a)+(a')\ \ge\ \frac{\kappa(B)}{\sigma},\qquad \kappa(B)=\min\Bigl(0.055,\ \frac{0.459}{4B+1}\Bigr).
\end{equation}
\end{proposition}

\begin{lemma}[positivity and periodicity] \proofst
\label{lem:P}
Let $p=\mathcal N_V\,\Gamma\,h/Z$ with $h$ as above, $\sum_nq0_n\le\Qb=\tfrac14$, $s\ge\sigma$ and $2\pi\sigma/\sqrt x\le1$; the density $\Gamma$ of $w$ may be any positive differentiable density.
Let $N_1=\sum_nn\,q0_n$,
\begin{equation}
a=\frac{\sqrt x\,N_1}{1-\Qb}+\frac{\Qb e^{-1/2}}{1-\Qb},\qquad \theta=\tfrac15e^{-10\pi^2/\tau}.
\end{equation}
Then every state satisfies
\begin{equation}
N_\phi\ \ge\ c_P\,\sigma^2\Jvv,\qquad c_P=\frac{\sigma^2}{8x}\,\frac{\theta}{1+\theta}\Bigl((1-b)-\frac{2a\sqrt x}{\pi}\Bigr)_+^2 .
\end{equation}
\end{lemma}

Proposition~\ref{prop:BE} is Lemma~BE of~\cite{H} with its two restrictions ($B\le10^4$, $\ell\le6907.8$) removed; for $B\le10^4$ its constant $0.459/(4B+1)\ge1.147\times10^{-5}$ agrees with the constant $1.1\times10^{-5}$ of~\cite{H}.
Lemma~\ref{lem:P} is new.
It does not use (P0), the regime of $x$, or any bound on $\eta$.

\subsection{Proof of Theorem~\ref{thm:main} from the five statements}

\begin{proof}
By Eq.~\eqref{eq:ga} it suffices to bound $(a)+(a')$.
Fix a state and write $B=b/\sigma^2$.

\emph{Case $55.26\le\ell\le960.8$.}
If $b>\tfrac12$, Theorem~\ref{thm:res} applies.
Let $b\le\tfrac12$.
On this range of $\ell$, $50(8\ell+100)\le389\,320<2.5\times10^{10}$.
If $B\ge2.5\times10^{10}$, Theorem~\ref{thm:narrow} gives $(a)+(a')\ge\Gamma_\ell$.
If $B<2.5\times10^{10}$, then $\tau=x/b\ge5\times10^{11}/B>20$, and item~(i) of Theorem~H1 of~\cite{H} gives $g\sqrt\ell\ge1-\bar\varepsilon$.

\emph{Case $\ell\ge960.8$.}
If $b>\tfrac12$, Theorem~\ref{thm:res} applies.
If $B\le8\ell+100$, Theorem~\ref{thm:verywide} applies.
If $8\ell+100\le B<\ell^4$, Theorem~\ref{thm:wide} applies.
If $B\ge\ell^4$ and $b\le\tfrac12$, then $B\ge\max(2.5\times10^{10},50(8\ell+100))$, because $\ell^4\ge8.5\times10^{11}$ and $\ell^4\ge50(8\ell+100)$ for $\ell\ge960.8$, and Theorem~\ref{thm:narrow} applies.
These cases exhaust the states, and each gives $(a)+(a')\ge\Gamma_\ell$, i.e.\ $g\sqrt\ell\ge(1-\bar\varepsilon)\sqrt{1+10/\ell}$.
\end{proof}

\subsection{Proof of Theorem~\ref{thm:verywide} from Proposition~\ref{prop:BE} and Lemma~\ref{lem:P}}

\begin{proof}
Let $\tau\ge\tau_0(\ell)$.
Proposition~\ref{prop:BE} and $B\le8\ell+100$ give $(a)+(a')\ge0.459\,e^{\ell}/(32\ell+401)$, which exceeds $\Gamma_\ell$ by a factor larger than $e^{950}$ for $\ell\ge960.8$.

Let $\tau<\tau_0(\ell)$.
Since $2A_w\ge\sigma^2\Jww$ (because $0\le\phi\le2$) and $\mathrm{den}=\sigma^2\Jvv+\sigma^2\Jww$, Lemma~\ref{lem:P} gives
\begin{equation}
2A_w+N_\phi\ \ge\ \sigma^2\Jww+c_P\,\sigma^2\Jvv\ \ge\ \min(1,c_P)\,\mathrm{den},
\end{equation}
so it suffices to show $c_P\ge\sigma\Gamma_\ell$.
The regime gives $\tau B=x/\sigma^2\ge5\times10^{11}$, hence $\tau\ge5\times10^{11}/(8\ell+100)$ and $10\pi^2/\tau\le10\pi^2(8\ell+100)/(5\times10^{11})$; also $x=\tau B\sigma^2<\tau_0(\ell)(8\ell+100)\sigma^2$, and the condition $2\pi\sigma/\sqrt x\le1$ of the lemma holds.
With $\theta/(1+\theta)\ge\tfrac16e^{-10\pi^2/\tau}$ and the bracket of $c_P$ above $1-10^{-400}$,
\begin{equation}
\begin{split}
\ln\frac{c_P}{\sigma\Gamma_\ell}\ \ge\ &\ell-\frac{10\pi^2(8\ell+100)}{5\times10^{11}}\\&-\ln\bigl(48\,\tau_0(\ell)(8\ell+100)\bigr)-\ln\Gamma_\ell-10^{-399}.
\end{split}
\end{equation}
The right side is at least $947$ at $\ell=960.8$ and increases with $\ell$.
\end{proof}

States with $\tau<\tau_0(\ell)$ exist in the region of Theorem~\ref{thm:verywide} only for $\ell>5.19\times10^{8}$ and $B>5\times10^{11}/\tau_0(\ell)$ (the hatched strip of Fig.~\ref{fig:map}); below that resolution Proposition~\ref{prop:BE} alone proves the theorem.

\subsection{What is new and what is borrowed}

\begin{table}
\caption{Origin of the statements used in Theorem~\ref{thm:main}.}
\label{tab:origin}
\begin{ruledtabular}
\small
\begin{tabular}{p{0.56\columnwidth}p{0.36\columnwidth}}
statement & origin\\ \hline
$g\ge(1-\bar\varepsilon)((a)+(a'))$, the conditions (P0), (P2)--(P4), (P6) & Paper~F, Eq.~(3), Table~5\\
period-average lemmas (Lemmas~19, 21--23 of F) & Paper~F, App.~D\\
Theorem~H1 ($\tau\ge20$, $\ell\le960.8$) & Paper~H, \Hsec{VIII}\\
Lemma~R$_\theta$ (local route), Lemma~CM-1, second line of Lemma~CN-1 & Paper~H, \Hsec{IX~A}\\
Lemma~A4*, the steps of Lemma~BE, window constants $U_1,U_2,\gamma_P,c_\chi$ & Paper~H, \Hsecs{III, VI, VII and Apps.~D, E}\\
the local lemmas (direct form, main-term identity, ring, pair; Lemmas~\ref{lem:lines}--\ref{lem:G0} and App.~\ref{app:A}) & new\\
Theorems~\ref{thm:narrow}, \ref{thm:res}, \ref{thm:wide}, \ref{thm:verywide} & new\\
Proposition~\ref{prop:BE} (Lemma~BE for all $B\le8\ell+100$, all $\ell\ge960.8$) & new; same core as Lemma~BE\\
Lemma~\ref{lem:P} (positivity and periodicity) & new\\
the brightness hypothesis (H$'$) in closed form, the lines and their constants & new (Secs.~\ref{sec:setting}--\ref{sec:wide})\\
\end{tabular}
\end{ruledtabular}
\end{table}

Table~\ref{tab:origin} separates the two.
Three remarks on the hypotheses.
First, $\eta\le1/\ell$ enters Theorems~\ref{thm:narrow}, \ref{thm:res} and~\ref{thm:wide} through the sandwich factors $1/(1\pm\eta)$ written into the lines, and Theorem~H1 through its item~(i); Theorem~\ref{thm:verywide} does not need it.
Second, the bound $q0_n\le2e^{-n^2x/2}$ used in Proposition~\ref{prop:BE} holds for every window of~\cite{F}, as noted in~\cite{H}.
Third, the factor $1-\bar\varepsilon$ is the hypothesis (P6); every statement of this paper is a statement about $(a)+(a')$ and is independent of the value of $\bar\varepsilon$.
